% ---
% titre: Droite numérique (placer, comparer, encadrer)
% theme: NO
% chapitre: Nombres rationnels
% sous_chapitre: Placer, comparer et encadrer
% niveau: [10e]
% type: EX
% tags: [c-droite-numerique, c-encadrement, c-comparaison, f-individuel, p-diagnostic]
% difficulte: 1
% source: latex
% ---

% Droite graduée de 16 cm : #1 = nombre à gauche, #2 = nombre à droite,
% #3 = nombre d'intervalles, #4 = points de la correction (visibles avec \sol)
\newcommand{\droite}[4]{%
  \begin{tikzpicture}
    \draw[->, >=stealth] (-0.4,0) -- (16.5,0);
    \foreach \i in {0,...,#3} {\draw ({\i*16/#3},0.15) -- ({\i*16/#3},-0.15);}
    \node[below] at (0,-0.2) {#1};
    \node[below] at (16,-0.2) {#2};
    #4
  \end{tikzpicture}%
}
% Point de la correction : abscisse en cm, nom du point
\newcommand{\pt}[2]{\draw[thick] (#1,0.3) -- (#1,-0.3); \node[above] at (#1,0.3) {\textbf{#2}};}

\header{Droite numérique : placer, comparer, encadrer}


\def\x{1.3cm}
\def\arraystretch{2.2}
\vfill
{\bfseries\color{themeColor}Exercice 1.} Les nombres entiers

\textbf{a)} Place sur la droite les points $A = 3$, $B = 8$ et $C = 5$.

\begin{center}
\droite{0}{10}{10}{\sol{}{\pt{4.8}{A}\pt{12.8}{B}\pt{8}{C}}}
\end{center}

\textbf{b)} Complète avec $<$ ou $>$ : \quad 7 \sol{\bl[\x]}{$>$} 4

\textbf{c)} Encadre 6 par deux nombres entiers qui se suivent : \quad \sol{\bl[\x]}{5} $<$ 6 $<$ \sol{\bl[\x]}{7}

\vfill
{\bfseries\color{themeColor}Exercice 2.}  Les nombres à virgule

\textbf{a)} Place sur la droite les points $A = 2{,}5$, $B = 4{,}1$ et $C = 0{,}8$.

\begin{center}
\droite{0}{5}{5}{\sol{}{\pt{8}{A}\pt{13.12}{B}\pt{2.56}{C}}}
\end{center}

\textbf{b)} Encadre 2,5 par deux nombres entiers qui se suivent : \quad \sol{\bl[\x]}{2} $<$ 2,5 $<$ \sol{\bl[\x]}{3}

\textbf{c)} Lequel des points $A$, $B$ et $C$ est le plus proche de 0 ? \quad \sol{\bl[\x]}{$C$}

\vfill
{\bfseries\color{themeColor}Exercice 3.} On zoome.

La droite va de 2 à 3. Il y a 10 intervalles égaux entre 2 et 3.

\textbf{a)} Place les points $A = 2{,}3$ et $B = 2{,}8$.

\begin{center}
\droite{2}{3}{10}{\sol{}{\pt{4.8}{A}\pt{12.8}{B}}}
\end{center}

\textbf{b)} Place le point $C = 2{,}25$. Entre quelles graduations se trouve-t-il ?

\sol{\bl[\x]}{2,2} $<$ 2,25 $<$ \sol{\bl[\x]}{2,3}

\begin{center}
\droite{2}{3}{10}{\sol{}{\pt{4}{C}}}
\end{center}

\textbf{c)} Encadre 2,25 par deux nombres entiers qui se suivent : \quad \sol{\bl[\x]}{2} $<$ 2,25 $<$ \sol{\bl[\x]}{3}

\newpage
{\bfseries\color{themeColor}Exercice 4.}  comparer et ranger

\textbf{a)} Complète avec $<$, $>$ ou $=$ :

\begin{tabularx}{\linewidth}{X c l @{\hspace{15pt}}|@{\hspace{15pt}} X c l @{\hspace{15pt}}|@{\hspace{15pt}} X c l}
7 & \sol{\bl[\x]}{$>$} & 4 & 4,9 & \sol{\bl[\x]}{$<$} & 5 & 0,5 & \sol{\bl[\x]}{$>$} & 0,45 \\
3,07 & \sol{\bl[\x]}{$<$} & 3,7 & 2,60 & \sol{\bl[\x]}{$=$} & 2,6 & 0,3 & \sol{\bl[\x]}{$>$} & 0,03 \\
\end{tabularx}

\def\arraystretch{1}
\medskip
\textbf{b)} Range du plus petit au plus grand : \quad 0,2 \; ; \; 0,25 \; ; \; 0,5 \; ; \; 0,05

\medskip
\sol{\bl[\x]}{0,05} $<$ \sol{\bl[\x]}{0,2} $<$ \sol{\bl[\x]}{0,25} $<$ \sol{\bl[\x]}{0,5}

\vspace{20pt}
{\bfseries\color{themeColor}Exercice 5.}  défi

Cet exercice est difficile. Essaie quand même : tout ce que tu sais m'intéresse, même si ce n'est pas complet.

\textbf{a)} Place les points $A = 0{,}75$, $B = \frac{1}{2}$, $C = 1{,}25$ et $D = \frac{3}{2}$.

\begin{center}
\droite{0}{2}{20}{\sol{}{\pt{6}{A}\pt{4}{B}\pt{10}{C}\pt{12}{D}}}
\end{center}

\textbf{b)} Range du plus petit au plus grand : \quad 0,75 \; ; \; $\frac{1}{2}$ \; ; \; 1,25 \; ; \; $\frac{1}{4}$

\medskip
\sol{\bl[\x]}{$\frac{1}{4}$} $<$ \sol{\bl[\x]}{$\frac{1}{2}$} $<$ \sol{\bl[\x]}{0,75} $<$ \sol{\bl[\x]}{1,25}

\medskip
\textbf{c)} Écris ce que tu sais déjà sur les fractions :

\bl[\linewidth]

\bl[\linewidth]
